% Section 8
\section{Numerical evidence and an application}\label{sec:evidence}

The results below are evidence, not proof. They check that the theory behaves as
stated in designs where the truth is known, and they show where the proved
intervals are informative. All intervals are nominal 95\%.

\subsection{Verification}

Every closed form, inclusion and counterexample in Sections~\ref{sec:rxc}--\ref{sec:inference}
has a numerical or symbolic check against brute force or an LP solver. Examples:
the homogeneous form reproduces JKST's bounds to $10^{-16}$; the two-row closed
form agrees with the LP in 300 random cases, and its district version
(Proposition~\ref{prop:D3}) in 126 cases plus 195 with pure vertices; the face test
of Proposition~\ref{prop:R2} agrees with brute force in 800 of 800 cases; and every
counterexample quoted in the readings (the singleton hull, the null-direction
projection failure, the changing-system limit $\min\{Z,1-Z\}$, the
district-reversal example, the three-point bracket design) is reproduced exactly.
Two parts are not machine-checked: the measure-theoretic steps of
Theorems~\ref{thm:sharp} and~\ref{thm:affine}, which rest on cited theorems, and the
remark on randomising over integer vertices, which is only a sketch and is not used.

\subsection{Monte Carlo design}

Unit rates are Beta (for $2\times2$) or Dirichlet (for $3\times3$) draws whose
conditional means are exactly linear in composition given the district. So M1
always holds, and M0 holds only without district shifts. Each design has 1000
replications. Coverage is the share of replications whose interval contains the
truth; an empty interval (an empty inflated set, or crossing ends) counts as not
covering. A coverage of 1 means every replication covered, and the exact one-sided
95\% lower confidence bound is then $0.05^{1/1000}=0.997$. Monte Carlo standard
errors are in parentheses. Widths are medians after clipping to $[0,1]$. The
projection interval for $B$ takes its Bhatia--Davis allowance from a projection
interval for $\bar m$ on the rung alone at level 0.005, with 0.04 on the
estimation part (Proposition~\ref{prop:bd}); the CLR interval for $B$ uses the
Bhatia--Davis allowance with a CLR input. Designs:
\begin{itemize}
\item A: $2\times2$, M1, 10 districts of 93 units (regime R1) and 93 districts of 30;
\item C: $2\times2$ pooled, with an implied rate close to 0 at one vertex, so
Slater's condition holds, but only just;
\item E: $2\times2$ pooled, compositions spanning $[0.02,0.98]$, with group $b$'s
mean rate exactly 0 at the right vertex and group $w$'s exactly 0 at the left, so
the identified set of the model-implied rate is a single point and Slater fails;
\item B: $3\times3$ pooled, four targets.
\end{itemize}

\begin{table}[t]
\centering\footnotesize
\setlength{\tabcolsep}{3pt}
\begin{tabular}{@{}lrc ccc ccc@{}}
\toprule
& & & \multicolumn{3}{c}{Model-implied rate $\Bt$} & \multicolumn{3}{c}{Actual rate $B$} \\
\cmidrule(lr){4-6}\cmidrule(lr){7-9}
Design & $q$ & DD & Conservative & CLR & Projection & Conservative & CLR & Projection \\
\midrule
A, $10\times93$ & 21 & 0.574 & 1 / 0.596 & 1 / 0.432 & 1 / 0.554 & 1 / 0.696 & 1 / 0.538 & 1 / 0.662 \\
A, $93\times30$ & 187 & 0.567 & 1 / 0.633 & 1 / 0.414 & 1 / 0.932 & 1 / 0.694 & 1 / 0.477 & 1 / 0.964 \\
C, boundary & 3 & 0.325 & 0.996 / 0.041 & 0.994 / 0.039 & 0.999 / 0.041 & 1 / 0.087 & 0.999 / 0.057 & 1 / 0.059 \\
E, Slater fails & 3 & 0.271 & 0.990 / 0.179 & \textbf{0.944} / 0.143 & 0.995 / 0.194 & 1 / 0.249 & 0.998 / 0.207 & 1 / 0.251 \\
B, loyalty & 12 & 0.835 & 1 / 0.592 & 1 / 0.557 & 1 / 0.581 & 1 / 0.675 & 1 / 0.659 & 1 / 0.679 \\
B, defection & 12 & 0.635 & 1 / 0.399 & 1 / 0.369 & 1 / 0.395 & 1 / 0.447 & 1 / 0.431 & 1 / 0.445 \\
B, demobilisation & 12 & 0.533 & 1 / 0.333 & 1 / 0.304 & 1 / 0.329 & 1 / 0.381 & 1 / 0.363 & 1 / 0.378 \\
B, mobilisation & 12 & 0.884 & 1 / 0.619 & 1 / 0.559 & 1 / 0.595 & 1 / 0.683 & 1 / 0.646 & 1 / 0.661 \\
\bottomrule
\end{tabular}
\caption{Monte Carlo, model holds: coverage / median width (1000 replications). $q=\dim\beta$; DD is the median Duncan--Davis width. Design E, CLR for $\Bt$: Monte Carlo standard error 0.007.}\label{tab:mc1}
\end{table}

\begin{table}[t]
\centering\footnotesize
\setlength{\tabcolsep}{4pt}
\begin{tabular}{@{}rrccccccc@{}}
\toprule
& & & Plug-in & \multicolumn{3}{c}{Model-implied rate $\Bt$} & \multicolumn{2}{c}{Actual rate $B$} \\
\cmidrule(lr){5-7}\cmidrule(lr){8-9}
Districts & $q$ & DD & cover / empty & CLR & Projection & Theorem~\ref{thm:finite} & CLR & Conservative \\
\midrule
30 & 61 & 0.571 & 0.305 / 0.241 & 0.999 / 0.557 & 1 / 0.967 & 1 / 1.000 & 1 / 0.746 & 1 / 1.000 \\
93 & 187 & 0.566 & 0.042 / 0.463 & 0.995 / 0.467 & 1 / 1.000 & 1 / 1.000 & 0.998 / 0.575 & 1 / 0.930 \\
200 & 401 & 0.566 & 0.005 / 0.717 & 0.996 / 0.416 & 1 / 1.000 & 1 / 1.000 & 1 / 0.491 & 1 / 0.890 \\
\bottomrule
\end{tabular}
\caption{Monte Carlo, many small districts (regime R2): M1, $2\times2$, 10 units per district. Coverage / median width. The projection interval and Theorem~\ref{thm:finite} for $B$ have coverage 1 and width 1.000 in every row.}\label{tab:mc2}
\end{table}

\subsection{Monte Carlo results}

Tables~\ref{tab:mc1} and~\ref{tab:mc2} report the designs in which the model holds.

\begin{enumerate}
\item \emph{No proved interval under-covers where its theorem applies.} The
projection interval covers at least 0.995 (design E, standard error 0.002) and the
conservative interval at least 0.990 (design E, 0.003). Theorem~\ref{thm:finite}
covered in every replication of every design where it was computed, including
regime R2 and designs A. In the R2 rows ($93\times30$ and 10 units per district),
$q$ grows with the number of districts, so the coverage of the projection and
conservative intervals there is itself a simulation finding.
\item \emph{CLR (a simulation finding) covers at least 0.994 in every design where
Slater's condition holds}, including R2 with 200 districts. At the Slater failure
(design E), CLR for $\Bt$ covers 0.944 (standard error 0.007), within one standard
error of nominal. This is the informative case: the identified set is a point, the
interval has no slack, and coverage falls from near one to the nominal level. The
coverage of 1 elsewhere reflects truths in the interior of the identified set.
\item \emph{The plug-in fails as Theorem~\ref{thm:limit} and
Proposition~\ref{prop:triangular} predict.} In R2 its coverage falls to 0.305, 0.042
and 0.005 at 30, 93 and 200 districts, and it is empty 24\%, 46\% and 72\% of the
time. In design E it covers 0.267 and is empty half the time: with a point
identified set, any sampling noise empties it.
\item \emph{The projection interval is informative only when $q$ is small.} Its
critical value grows like $\sqrt q$: 2.8 at $q=3$, 4.6 at $q=12$, 14.8 at $q=187$ and
21.2 at $q=401$, against about 2--3 for CLR. In pooled designs it is only 0.01--0.03
wider than CLR; with 10 districts it is 0.554 against 0.432; with 93 districts it
is 0.93, and $[0,1]$ at 10 units per district. In grouped designs with many districts,
CLR is the only informative interval, and its support there is simulation only.
\item \emph{Theorem~\ref{thm:finite} is valid but uninformative}: clipped width
1.000 in every M1 design, including $10\times93$, where CLR is 0.43 and
Duncan--Davis 0.57.
\item \emph{In R2, CLR narrows as districts are added}: for $\Bt$, 0.557, 0.467 and
0.416 at 30, 93 and 200 districts. With 200 districts of 10 units it is narrower
than Duncan--Davis (0.566).
\item \emph{The residual term is a large part of the width.} Targeting $\Bt$ instead
of $B$ narrows CLR by 0.06--0.10 in the $3\times3$ design and 0.06--0.19 in the
$2\times2$ M1 designs. The Bhatia--Davis allowance narrows CLR for $B$ from 0.086 to
0.057 in the boundary design C; where the plausible rates include 0.5 it equals
Popoviciu's and is up to 0.015 wider, because part of $\alpha$ is spent on $\bar m$.
The projection interval for $B$ with the proved allowance of
Proposition~\ref{prop:bd} differs in width by at most 0.002 from the version with a
CLR input, and its coverage is the same in every design.
\item \emph{Rank.} The numerical rank of $\hat\Sigma$ equals $\dim\beta$ in every
replication of every $2\times2$ design ($\dim\beta$ from 3 to 401). The $3\times3$
fit on all three columns has rank 12 of 18 because of the exact adding-up identity;
on the $C-1$ free columns it has full rank 12.
\end{enumerate}

\paragraph{District shocks.} Table~\ref{tab:shocks} compares M0 and M1 in a
$2\times2$ design with 20 districts of 50 units. The district shifts are added to
both groups' rates and depend on the district's mean composition $m_d$: none (S0);
affine in $m_d$ plus noise (S1); or a common shock $0.25\cos(2\pi m_d)$, which is not
affine (S2). With a non-affine shock every M0 interval fails, mostly by being empty,
while M1 covers in every replication and the pooling test always rejects. With
affine shifts the pooling test rejects every time, because M0 is false for the
observed design (Proposition~\ref{prop:tests}(c)), yet M0's intervals still cover:
shifts affine in composition change the pooled conditional mean almost only
through terms the pooled model absorbs (Proposition~\ref{prop:D4}(b)). A rejection
of the pooling test says that M0 is wrong, not that its intervals are.

\begin{table}[t]
\centering\footnotesize
\setlength{\tabcolsep}{3pt}
\begin{tabular}{@{}llccccc@{}}
\toprule
Shifts & Fit & Plug-in & CLR, $\Bt$ & Projection, $\Bt$ & Projection, $B$ & Pooling test \\
\midrule
S0 none & M0 & 1 / 0.000 & 1 / 0.329 & 1 / 0.336 & 1 / 0.436 & \\
 & M1 & 0.999 / 0.000 & 1 / 0.468 & 1 / 0.704 & 1 / 0.808 & size 0.090 ($\chi^2$) \\
S1 affine & M0 & 1 / 0.000 & 1 / 0.335 & 1 / 0.344 & 1 / 0.444 & \\
 & M1 & 0.938 / 0.000 & 1 / 0.437 & 1 / 0.672 & 1 / 0.773 & rejects 1.000 \\
S2 non-affine & M0 & 0.002 / 0.987 & 0.009 (0.963) & 0.013 (0.950) & 0.206 (0.478) & \\
 & M1 & 0.997 / 0.002 & 1 / 0.328 & 1 / 0.509 & 1 / 0.611 & rejects 1.000 \\
\bottomrule
\end{tabular}
\caption{District shocks, M0 against M1: $2\times2$, 20 districts $\times$ 50 units, 1000 replications. Coverage / median width; for S2 under M0, coverage with the share of empty intervals in parentheses. The plug-in column gives coverage / share empty. The pooling-test column gives the $\chi^2$ rejection rate at 5\%.}\label{tab:shocks}
\end{table}

\paragraph{Calibration of the tests.} In design S0 the pooling test has 38
restrictions. Its HC3 Wald statistic has mean 39.7 and variance 100, against 38 and
76 for $\chi^2_{38}$, and the $\chi^2$ test rejects 0.090 (standard error 0.009) at
nominal 5\% in the 1000-replication study above. A separate 2000-replication run
with the restricted wild bootstrap ($B=399$) gives:

\begin{center}\small
\begin{tabular}{@{}lccccc@{}}
\toprule
Test & Restrictions & Replications & $\chi^2$ & F & Wild bootstrap \\
\midrule
Pooling (M0 against M1) & 38 & 2000 & 0.078 (0.006) & 0.072 (0.006) & 0.043 (0.005) \\
Curvature (M1 against M2) & 19 & 1000 & 0.057 (0.007) & 0.055 (0.007) & 0.054 (0.007) \\
\bottomrule
\end{tabular}
\end{center}

The two $\chi^2$ figures for the pooling test (0.090 and 0.078) come from different
runs and differ by about one combined standard error. The bootstrap is close to
nominal in this design; that is a simulation finding.

\subsection{A benchmark with known truth: New Zealand elections}

New Zealand publishes, for each electorate, the exact cross-tabulation of the party
vote by the electorate (candidate) vote, together with polling-place totals of both.
This gives a benchmark in which the truth is known. Units are polling places; rows
are the party vote (Labour, National, other); columns are the electorate vote
(Labour candidate, National candidate, other); seven general elections, 2002--2020.
Intervals are for the actual rate, use the Bhatia--Davis allowance, and are
intersected with Duncan--Davis. The benchmark was run with version 0.2.0 of the
package, before the later changes to the tests and to the variance bound of the
projection interval for $B$ (Proposition~\ref{prop:bd}); those changes can move the
projection and test columns slightly.

\paragraph{One table per electorate.} There are 492 electorate tables with a median
of 64 polling places each, and 4302 cells. The pooled model's CLR and projection
intervals covered every cell, at about Duncan--Davis width (median width ratio
1.00; CLR median width 0.833 against 0.843 for DD). Adding the ordering belief that
each party's voters are at least as likely as the other party's voters to vote for
that party's candidate also covered every cell. Goodman's constancy (L6) gave
intervals with median width 0.415 that covered 0.939 of cells and missed the truth
with a non-empty interval in 2.2\%. The plug-in estimate of the M0 identified set was
empty in 48\% of cells. With about 64 units per table, the valid intervals are about
as wide as Duncan--Davis.

\paragraph{Electorates as districts.} Pooling each election's 62--69 electorates
(3398--6550 polling places) and fitting M1 with electorates as districts, the data
reject the model: the M1 plug-in estimate is empty in all seven elections, and the
curvature and pooling tests have $p\approx0$ (the pooling test has 366--408
restrictions). The emptiness test itself rejects the pooled model in two elections
and M1 in none. Over the 63 cells (seven elections, nine cells), the M1 CLR
intervals had crossing ends (or ends outside the logical range) in 21\% of cells and missed the truth with a non-empty
interval in 8\% (5 of 63); the projection interval never missed, but it was as wide
as Duncan--Davis. Where the model is wrong, the tests flag it, and a non-empty CLR
interval is not reassurance.

\subsection{The sign module}

\paragraph{How often signs add information.} In 2000 random $2\times2$ LCE designs
with true signs imposed, the within sign strictly tightened JKST's interval in 83\%
of designs (median width ratio 0.70), the between sign in 89\% (0.67), and CR, with
unknown sign, in 80\% of the 1060 designs where CR holds (0.69). In a $3\times2$
design with true-sign rows on both associations, a row was implied in 8\% of
target-row pairs, cut the set but not the target in 28\%, and tightened the target in
64\% (Proposition~\ref{prop:R2}).

\paragraph{CR at the edge.} Table~\ref{tab:cr} uses a $2\times2$ design with 150
units, 600 replications, Beta rates with linear means (so LCE and CR hold) and the
conservative interval for the gap $\tilde D$. In the first row the truth lies on the
edge of the CR band, so $\delta^*_W=0$.

\begin{table}[t]
\centering\footnotesize
\begin{tabular}{@{}lcccc@{}}
\toprule
Design & JKST alone & CR slab & CR branch union & CR plug-in bracket \\
\midrule
Truth on the CR edge, $\ER^*=0.15$ & 1.000 / 1.106 & 0.992 / 0.359 & 0.992 / 0.357 & \textbf{0.917} / 0.256 \\
Truth on the edge, $\ER^*=0.005$ & 1.000 / 1.223 & 1.000 / 0.242 & 1.000 / 0.242 & 1.000 / 0.147 \\
Interior truth & 1.000 / 1.232 & 1.000 / 0.480 & 1.000 / 0.480 & 1.000 / 0.388 \\
\bottomrule
\end{tabular}
\caption{Contextual reinforcement in $2\times2$: coverage / median length of conservative intervals for the gap (600 replications). The branch union is measured by its occupied length. No interval was empty.}\label{tab:cr}
\end{table}

The slab cuts the median length from 1.106 to 0.359 at the edge and keeps coverage.
The plug-in bracket, which imposes $\hull\{\gamma\widehat\ER,\widehat\ER\}$ and then
adds the target's own inflation, covers 0.917, consistent with the limit of
Theorem~\ref{thm:R10}(e) (at most 0.9335 in the three-point design). The literal
intersection of the fitted bracket with a finished JKST interval covered 304 of 600
(51\%), in line with its limit of at most one half. The package's default CLR branch
union, in a separate run of 200 replications of the edge design, covered 0.990 with
median width 0.230; that interval is algorithmic.

\paragraph{A false belief.} With the same generator and means $m_0=.40-sx$ and
$m_1=.50+s-sx$, the aggregate regression is $.40+.10x$ for every $s$, the between
sign is positive throughout, and for $s>0$ both contextual slopes are negative, so
CR is false.

\begin{center}\small
\begin{tabular}{@{}rrccccc@{}}
\toprule
Units & $s$ & True $\tilde D$ & JKST covers & CR slab covers & JKST width & CR width \\
\midrule
150 & 0 & .100 & 400/400 & 399/400 & 1.262 & .306 \\
150 & .08 & .170 & 400/400 & 336/400 & 1.256 & .304 \\
150 & .25 & .318 & 400/400 & 2/400 & 1.254 & .303 \\
600 & 0 & .100 & 400/400 & 398/400 & 1.074 & .194 \\
600 & .08 & .169 & 400/400 & 110/400 & 1.069 & .196 \\
600 & .25 & .317 & 400/400 & 0/400 & 1.066 & .195 \\
\bottomrule
\end{tabular}
\end{center}

No interval was empty. A false restriction gives narrow, non-empty and reproducibly
wrong intervals: precision makes the wrong answer more reproducible, not more
credible. This is the price stated in Theorem~\ref{thm:R10}(d).

\subsection{A worked example}\label{sec:example}

We illustrate the package on the \code{voter\_vax} data distributed with the
MonotoneEI package of EGGHM (and included in \pkg{eiPartialID}'s tests): 2059 US
counties, with $X$ the Republican share of the presidential vote, $T$ the share
vaccinated against COVID-19, and $N$ the total votes cast. Rows are Republican and
other voters; the column is ``vaccinated''; the target is the gap $D$ between the two
groups' vaccination rates. The example illustrates the mechanics. It is not a
substantive finding: the vaccination share refers to the county population and the
group shares to voters, so the closed-population assumption A1 holds only
approximately.

Fitting the pooled rung with size weights,
\begin{center}
\code{eiBounds(x = prop\_republican, t = prop\_covid\_vax, n = total\_votes, weights = "n",}\\
\code{targets = list(D = matrix(c(1, -1, 0, 0), 2, 2)), ci = "clr")}
\end{center}
gives the results in Table~\ref{tab:vv} (seed 1). The estimated identified set for
the gap under LCE alone is $[-0.486,-0.001]$: Republican voters' vaccination rate is
lower, by between 0 and 49 points. The projection interval, which carries the
asymptotic guarantee for the pooled model, is $[-0.607,0.080]$, and the CLR interval
is slightly narrower. The actual-rate intervals are much wider, because the residual
allowance covers each county's local luck.

\begin{table}[t]
\centering\small
\begin{tabular}{@{}lcccc@{}}
\toprule
& \multicolumn{2}{c}{M0 alone} & \multicolumn{2}{c}{M0 + CR} \\
\cmidrule(lr){2-3}\cmidrule(lr){4-5}
& Gap $D$ & Republican rate & Gap $D$ & Republican rate \\
\midrule
Estimate (identified set) & $[-0.486,-0.001]$ & $[0.345,0.597]$ & $[-0.470,-0.050]$ & $[0.354,0.571]$ \\
CLR, $\Bt$ (algorithmic) & $[-0.571,\ \ 0.056]$ & $[0.301,0.623]$ & $[-0.525,-0.044]$ & $[0.330,0.580]$ \\
Projection, $\Bt$ & $[-0.607,\ \ 0.080]$ & $[0.282,0.634]$ & $[-0.542,-0.043]$ & $[0.323,0.582]$ \\
Conservative, $\Bt$ & $[-0.599,\ \ 0.112]$ & $[0.291,0.651]$ & $[-0.530,\ \ 0.011]$ & $[0.326,0.600]$ \\
CLR, $B$ (algorithmic) & $[-0.750,\ \ 0.233]$ & $[0.229,0.694]$ & $[-0.702,\ \ 0.131]$ & $[0.259,0.650]$ \\
Projection, $B$ & $[-0.763,\ \ 0.257]$ & $[0.210,0.705]$ & $[-0.719,\ \ 0.133]$ & $[0.252,0.652]$ \\
Duncan--Davis & $[-0.763,\ \ 0.547]$ & $[0.202,0.881]$ & $[-0.763,\ \ 0.547]$ & $[0.202,0.881]$ \\
\bottomrule
\end{tabular}
\caption{The \code{voter\_vax} example: 2059 counties, pooled rung, size-weighted regression, 95\% intervals. The conservative row comes from a separate fit with \code{ci = "conservative"}.}\label{tab:vv}
\end{table}

\paragraph{Descriptive quantities.} The Republican share has mean $\mu=0.482$ and
between-county variance $0.0267$, so $\gamma=0.107$. The regression slope is
$\ER=-0.470$ and the neighbourhood gap is $\NM=\gamma\ER=-0.050$. By
Proposition~\ref{prop:G2}, each sign belief cuts the identified set at one of these
two numbers:
\begin{center}\small
\begin{tabular}{@{}ll@{}}
\toprule
Belief & Estimated set for $D$ \\
\midrule
$\delta_W\ge0$ (vaccination rises with the Republican share, on average) & $[-0.486,-0.470]$ \\
$\delta_W\le0$ & $[-0.470,-0.001]$ \\
$\delta_B\ge0$ (Republicans more likely to be vaccinated within a county) & $[-0.050,-0.001]$ \\
$\delta_B\le0$ & $[-0.486,-0.050]$ \\
CR (the two associations share a sign) & $[-0.470,-0.050]$ \\
\bottomrule
\end{tabular}
\end{center}
Under the rung alone neither association's sign is settled: the estimated set for
$\delta_B$ is $[-0.488,0.055]$ and for $\delta_W$ is $[-0.525,0.018]$.
\code{eiRedundancy()} reports that CR tightens the estimated set for $D$ by 0.066,
whether added first or last. Since $\ER<0$, only the negative branch of CR is
non-empty (Proposition~\ref{prop:tworowsign}(d)), and the CR set is the band between
$\ER$ and $\NM$. With CR, the projection interval for $\Bt$, $[-0.542,-0.043]$,
excludes zero: under LCE and CR, Republican voters' model-implied vaccination rate
is lower, with an asymptotic guarantee for the pooled model
(Theorem~\ref{thm:R10}). The conclusion is only as credible as CR.

\paragraph{A joint rejection.} Imposing both $\delta_W\ge0$ and $\delta_B\ge0$ asks
for $D\le\ER=-0.470$ and $D\ge\NM=-0.050$ at once, which is impossible. The estimated
set is empty, and the emptiness test rejects the combination of the rung and the two
beliefs, while the rung alone is not rejected. The package reports exactly this: the
empty estimate is a sample diagnostic, the emptiness test is the test, the rejection
is of the combination and does not say which part fails, and each belief on its own
gives a non-empty set (Theorem~\ref{thm:R3}(e)). The signed associations have
effective numbers of counties of 271 ($\delta_B$) and 188 ($\delta_W$), so no single
county dominates them.
